% intro.tex (국문 v15.3, 2026-10-22): 영문 v15.3 의 쌍둥이. 네 문단: 문제(관측의 편중), 두 모델 유형과 제품의 공백, 접근, 기여.
활동층은 영구동토 위에 놓여 여름에 녹고 겨울에 다시 어는 지표 부근의 토층이다. 활동층 두께(ALT)는 계절 최대 융해 깊이이며, 지반 안정성, 지표 수문, 분해에 노출되는 토양 유기탄소의 양을 좌우한다\cite{Hjort2018,Schuur2015,Biskaborn2019}. ALT 는 해마다 열 조건에 따라 달라지며, 이 연도 간 변동은 환북극 활동층 관측망(CALM) 같은 장기 관측 프로그램이 기록한다\cite{Nelson1997,Streletskiy2025}. 그러나 관측은 드물고 군집되어 있다. 대부분은 소수의 캠페인과 장기 격자에서 나오고, 많은 지점은 한 번만 방문되며, 영구동토 지역의 넓은 부분에는 관측이 없다\cite{Moore2025,Veremeeva2025}. 따라서 ALT 지도는 드문 관측을 관측이 없는 격자 셀과 연도로 확장하는 모델로 만들어지며, 지도의 신뢰도는 그 모델이 관측 밖에서 유지하는 정확도에 좌우된다.

이 확장에 쓰는 모델은 물리 모델과 기계학습 모델로 나뉜다. Stefan 해는 토양의 열적 성질과 수분을 요약하는 계수 하나로 융해 깊이를 그 계절의 융해 도일(TDD)에 연결한다\cite{Stefan1891,Riseborough2008,Nelson1997}. Stefan 모델은 연도별 기후에 반응하며, 매개변수가 하나이므로 소수의 관측으로 다시 맞추어 관측 지점에서 먼 지역에도 적용할 수 있다. 약점은 지역당 계수 하나로는 지역 안의 토양과 지형 차이를 나타낼 수 없어 체계적 잔차가 남는다는 것이다. 기계학습(ML)의 장단점은 이와 반대이다. ML 은 기후 · 지형 · 토양 예측 변수로 학습 지점의 관측을 잘 재현하지만, 기후에 대한 ALT 의 반응을 학습용 관측값(이하 라벨)에서 배워야 하며, 공간적으로 분리된 자료로 검증하면 학습 지점과 조건이 다른 위치에서 정확도가 낮아진다\cite{Roberts2017,Ploton2020,MeyerPebesma2022,Wadoux2021}. 현재 북반구 영구동토 지역을 포괄하는 격자 ALT 제품도 이 두 방식 중 하나를 따른다. 유럽우주국 기후변화 이니셔티브(ESA CCI) 제품은 위성 지표 온도로 구동한 지중 열전달 모델을 풀고\cite{Westermann2024,Westermann2025}, Liu 등\cite{Liu2024}과 Wei 등\cite{Wei2026}의 ML 제품은 지점별 · 연도별 관측(이하 지점-연도 관측)을 연도별 예측 변수에 연결하며, Aalto 등\cite{Aalto2018}과 Ran 등\cite{Ran2022a}의 제품은 다년 평균을 지도화한다. ML 제품은 각자의 지점 집합 안에서 무작위 분할이나 지점 제외 분할로 검증되었다. 검증 자료가 서로 달라 보고된 오차(21.6 cm\cite{Liu2024}--60.5 cm\cite{Wei2026})는 비교할 수 없고, 학습 지점이 없는 지역에서 특정 연도의 관측으로 검증된 제품은 없다.

물리 모델이 주된 반응을 나타내고 학습기가 물리 모델이 설명하지 못한 부분을 보정하는 물리 기반 학습은 호수 · 하천 · 토양 응용에서 정확도와 외삽 성능을 개선하였다\cite{Willard2022,Read2019,Jia2021}. 알래스카에서는 학습에서 제외한 지점에서 Stefan 모델이 랜덤 포레스트보다 정확도를 더 잘 유지하였다\cite{Gautam2025}. Stefan 모델을 보정하는 학습기가 학습에 쓰지 않은 지역에서의 정확도(이하 전이성)를 유지하면서 ML 수준의 정확도를 얻는지는 ALT 에 대해 시험되지 않았다. 이러한 모델의 실용성은 보정이 Stefan 모델을 개선하는 조건, 새 지역에 필요한 지역 관측의 수, 같은 지점 · 같은 연도에서 평가한 직접 ML(공변량만으로 ALT 를 학습하는 ML) 및 공개 제품과의 정확도 비교로 판단하며, 이 세 사항은 확인되지 않았다.

이 연구는 이 세 사항을 지점-연도 설계로 다룬다. 이 설계에서 라벨, 기후 입력, 예측, 평가는 모두 특정 연도의 한 관측 지점 단위로 정의하며, 기후 입력은 해당 연도의 ERA5-Land 자료를 관측일까지 누적하여 계산한다\cite{munozsabater2021_era5_land}. 알래스카, 레나 델타, 캐나다 북서부의 2010--2024년 지점-연도 관측 \nncyAll{}개로 세 방법, 즉 Stefan 모델, 직접 ML, 그리고 지역 라벨로 계수를 다시 맞춘(재보정) Stefan 모델의 잔차를 입력 23개로부터 gradient boosting 으로 학습하는 물리 잔차 워크플로를 비교한다. 비교에는 무작위 분할부터 한 지역 전체의 홀드아웃까지 다섯 검증 설계를 쓰며, 공개 제품 5종도 같은 지점-연도에서 평가한다. 이 연구는 (i) 검증 설계에 따른 방법 순위의 변화와 그 원인, 직접 ML 의 정확도가 크게 낮아지는 조건을 제시하고, (ii) 잔차 학습기가 Stefan 모델을 개선하는 조건과 필요한 라벨 수를 밝히며, (iii) 공개 제품의 지역 간 정확도 차이를 설명하고, (iv) 적용 가능 영역(공변량이 학습 자료의 범위 안에 있어 예측을 적용할 수 있는 범위)과 예측 구간을 포함한 2011--2024년 연도별 1 km ALT 지도를 세 지역에 대해 제공한다.
